\section{Part II.a (i): Geometry and the Vacuum Equation}

\subsection{Conformal}

A position-dependent scale field $\sigma$, the same in all directions, induces a conformal metric $g = e^{2\sigma}\delta$. The Riemannian geometry of such a metric is entirely determined by derivatives of $\sigma$: the Christoffel symbols are polynomial in $\nabla \sigma$, and curvature becomes the second-order inhomogeneity of the scale field. This file computes the Levi-Civita data and derives the master regression theorem stating that the Riemann tensor is the Kulkarni--Nomizu product of the bare Euclidean metric $\delta$ with a scale deformation tensor $D$ built from $\sigma$.

\begin{definition}
The scale gradient: $\sigma_i = \partial_i \sigma$. \lean{sig}
\end{definition}

\begin{definition}
The scale Hessian: $\sigma_{ij} = \partial_i \partial_j \sigma$. \lean{hess}
\end{definition}

\begin{theorem}
The scale Hessian is symmetric: $\sigma_{ij} = \sigma_{ji}$. \lean{hess\_symm}
\end{theorem}

\begin{proof}
By commutativity of mixed partial derivatives.
\end{proof}

\begin{theorem}
Mixed derivatives of the scale gradient commute: $\partial_i (\partial_j \sigma) = \partial_j (\partial_i \sigma)$. \lean{d\_sig\_comm}
\end{theorem}

\begin{proof}
Immediate from symmetry of the Hessian.
\end{proof}

\begin{definition}
The flat Laplacian of the scale: $\Delta\sigma = \sum_i \sigma_{ii}$. \lean{lap}
\end{definition}

\begin{definition}
The squared scale gradient: $|\nabla\sigma|^2 = \sum_i \sigma_i^2$. \lean{gradsq}
\end{definition}

\begin{theorem}
A constant factors out of the squared gradient sum: $\sum_m c(\sigma_m^2) = c \cdot |\nabla\sigma|^2$. \lean{sum\_mul\_gradsq}
\end{theorem}

\begin{proof}
Immediate by distributivity.
\end{proof}

\begin{theorem}
A constant summed over $\text{Fin}\,n$ yields $n$ times the constant: $\sum_{a=1}^n c = nc$. \lean{sum\_const\_fin}
\end{theorem}

\begin{proof}
Immediate from the definition of finite sums.
\end{proof}

The scale connection is the one-form $A_i = \sigma_i$. Its curvature $F_{ij} = \partial_i A_j - \partial_j A_i$ vanishes identically by commutativity of mixed partials, so the scale connection is integrable. This absence of a second clock effect is why an integrable scale geometry succeeds where Weyl's 1918 non-integrable unified theory failed.

\begin{definition}
The scale connection one-form: $A_i = \sigma_i$. \lean{scaleConn}
\end{definition}

\begin{definition}
The curvature of the scale connection: $F_{ij} = \partial_i A_j - \partial_j A_i$. \lean{scaleCurv}
\end{definition}

\begin{theorem}
The scale connection is flat: $F_{ij} = 0$. \lean{scaleCurv\_eq\_zero}
\end{theorem}

\begin{proof}
By commutativity of mixed partial derivatives, $\partial_i(\partial_j \sigma) = \partial_j(\partial_i \sigma)$, so the curvature vanishes identically.
\end{proof}

\begin{definition}
The Christoffel symbols of $g = e^{2\sigma}\delta$:
$$\Gamma^a_{bc} = \delta^a_b \sigma_c + \delta^a_c \sigma_b - \delta_{bc}\sigma^a,$$
where indices are raised and lowered with $\delta$.  The conformal factor cancels entirely, leaving a polynomial in $\nabla\sigma$; the connection sees only changes of scale, not the scale itself. \lean{Chr}
\end{definition}

\begin{theorem}
The Christoffel symbols are symmetric in their lower indices: $\Gamma^a_{bc} = \Gamma^a_{cb}$. \lean{Chr\_symm}
\end{theorem}

\begin{proof}
Immediate from the definition and symmetry of the Kronecker delta.
\end{proof}

\begin{definition}
The Riemann tensor: $R^a_{bce} = \partial_c \Gamma^a_{be} - \partial_e \Gamma^a_{bc} + \sum_m (\Gamma^a_{cm}\Gamma^m_{be} - \Gamma^a_{em}\Gamma^m_{bc})$. \lean{Rm}
\end{definition}

\begin{definition}
The Ricci tensor: $R_{be} = \sum_a R^a_{bae}$. \lean{Ric}
\end{definition}

\begin{definition}
The bare scalar curvature (the $\delta$-trace of Ricci): $R_{\text{bare}} = \sum_b R_{bb}$. The true scalar curvature of $g$ is $e^{-2\sigma}$ times this; the conformal factor is kept explicit so that no invertibility hypothesis is ever needed. \lean{RscBare}
\end{definition}

\begin{theorem}
The sum of derivatives of Christoffel symbols: $\sum_a \partial_a \Gamma^a_{be} = 2\sigma_{be} - \delta_{be}\Delta\sigma$. \lean{sum\_d\_Chr}
\end{theorem}

\begin{proof}
Expanding $\Gamma^a_{be}$ and taking derivatives yields three terms summed over $a$. After applying the product rule and regrouping by $i$-values, this simplifies to the stated formula.
\end{proof}

\begin{theorem}
The trace of Christoffel symbols: $\sum_a \Gamma^a_{am} = n\sigma_m$. \lean{sum\_Chr\_diag}
\end{theorem}

\begin{proof}
Evaluating $\Gamma^a_{am} = \delta^a_a \sigma_m = \sigma_m$ and summing over $a = 1, \ldots, n$ gives $n\sigma_m$.
\end{proof}

\begin{theorem}
The derivative of the trace: $\sum_a \partial_e \Gamma^a_{ba} = n\sigma_{eb}$. \lean{sum\_d\_Chr\_trace}
\end{theorem}

\begin{proof}
Use $\Gamma^a_{ba} = \Gamma^a_{ab}$ by symmetry, then sum $\partial_e \Gamma^a_{ab} = \partial_e (\Gamma^a_{aa}\sigma_b) = n\partial_e \sigma_b = n\sigma_{eb}$.
\end{proof}

\begin{theorem}
A weighted sum of Christoffel symbols: $\sum_m (n\sigma_m) \Gamma^m_{be} = n(2\sigma_b\sigma_e - \delta_{be}|\nabla\sigma|^2)$. \lean{sum\_ChrTrace\_Chr}
\end{theorem}

\begin{proof}
Expanding $\Gamma^m_{be}$ and distributing the sum yields three terms, each involving Kronecker deltas. Regrouping produces the stated form.
\end{proof}

\begin{theorem}
The master contraction of the Riemann tensor, in full generality:
\begin{align*}
\sum_m \Gamma^a_{cm}\Gamma^m_{be} &= 2\delta_{ac}\sigma_b\sigma_e - \delta_{ac}\delta_{be}|\nabla\sigma|^2 \\
&\quad + \delta_{ab}\sigma_c\sigma_e + \delta_{ae}\sigma_b\sigma_c \\
&\quad - \delta_{cb}\sigma_a\sigma_e - \delta_{ce}\sigma_a\sigma_b.
\end{align*}
Every later curvature statement is a specialization of this identity. \lean{sum\_Chr\_Chr\_full}
\end{theorem}

\begin{proof}
Expanding both Christoffel symbols and collecting terms by their Kronecker delta factors yields the six terms shown.
\end{proof}

\begin{theorem}
The doubly contracted case: $\sum_{a,m} \Gamma^a_{em}\Gamma^m_{ba} = (n+2)\sigma_b\sigma_e - 2\delta_{be}|\nabla\sigma|^2$. \lean{sum\_Chr\_Chr}
\end{theorem}

\begin{proof}
Apply the master contraction with $a \to a$ and $c \to e$, then sum over $a$. Using $\delta_{aa} = n$ and the definitions of Laplacian and squared gradient gives the result.
\end{proof}

\begin{theorem}
Ricci curvature is the second-order inhomogeneity of the scale:
$$R_{be} = -(n-2)(\sigma_{be} - \sigma_b\sigma_e) - \delta_{be}(\Delta\sigma + (n-2)|\nabla\sigma|^2).$$
Read physically: gravity is the failure of the local unit of measure to vary affinely across the substrate. \lean{Ric\_eq}
\end{theorem}

\begin{proof}
The Ricci tensor can be split into three parts: $\sum_a \partial_a \Gamma^a_{be}$, $\sum_a \partial_e \Gamma^a_{ba}$, and the quadratic terms. Using the contraction theorems, each part contributes to the final formula.
\end{proof}

\begin{theorem}
The scalar curvature of $g = e^{2\sigma}\delta$:
$$e^{2\sigma}R = -2(n-1)\Delta\sigma - (n-1)(n-2)|\nabla\sigma|^2.$$
For $n=3$ this is the Hamiltonian constraint of a conformally flat slice; its linearization is Poisson's equation. \lean{RscBare\_eq}
\end{theorem}

\begin{proof}
Sum the Ricci components $R_{bb}$ using the Ricci formula and collect terms involving the Laplacian and squared gradient.
\end{proof}

The entire Riemann tensor is a linear function of one symmetric tensor built from $\sigma$. This deformation tensor measures the obstruction to the local unit varying affinely.

\begin{definition}
The scale deformation tensor:
$$D_{ij} = 2\sigma_{ij} - 2\sigma_i\sigma_j + \delta_{ij}|\nabla\sigma|^2,$$
the obstruction to the local unit of measure varying affinely. (The factor of 2 is carried explicitly so that no division by 2 is required, keeping every statement valid over an arbitrary commutative ring.) \lean{Defm}
\end{definition}

\begin{theorem}
The deformation tensor is symmetric: $D_{ij} = D_{ji}$. \lean{Defm\_symm}
\end{theorem}

\begin{proof}
By symmetry of the Hessian, commutativity of multiplication, and symmetry of the Kronecker delta.
\end{proof}

\begin{theorem}
The derivative part of the Riemann tensor:
\begin{align*}
\partial_c \Gamma^a_{be} - \partial_e \Gamma^a_{bc} &= \delta_{ae}\sigma_{bc} - \delta_{be}\sigma_{ac} \\
&\quad - \delta_{ac}\sigma_{be} + \delta_{bc}\sigma_{ae}.
\end{align*}
\lean{d\_Chr\_diff}
\end{theorem}

\begin{proof}
Expand the Christoffel symbols and apply the product rule to each term, then use commutativity of mixed partials to simplify.
\end{proof}

\begin{theorem}
The master regression theorem: curvature is scale deformation. The Riemann tensor of $g = e^{2\sigma}\delta$ satisfies
$$2R^a_{bce} = \delta_{ae}D_{bc} - \delta_{ac}D_{be} + \delta_{bc}D_{ae} - \delta_{be}D_{ac}.$$
The Riemann tensor is exactly the Kulkarni--Nomizu product of the bare Euclidean bookkeeping $\delta$ with the scale deformation $D$. Nothing about curvature is independent of the scale field: geometry is bookkeeping times scale deformation, and nothing else. \lean{two\_Rm\_eq}
\end{theorem}

\begin{proof}
Split the Riemann tensor into derivative and quadratic parts. The derivative part is given by the preceding theorem. The quadratic part simplifies using the master contraction theorem. Combining these yields the Kulkarni--Nomizu form.
\end{proof}

Over any ring in which 2 is invertible---in particular, any $\mathbb{R}$-algebra---the master identity gives at once that a deformation-free scale field is exactly a flat geometry. Vacuum is a statement about the uniformity of scale change, never about the absence of scale.

\begin{theorem}
Deformation-free scale implies flat space. Not merely Ricci-flat: the whole Riemann tensor vanishes.
$$\forall i,j\,: D_{ij} = 0 \implies R^a_{bce} = 0.$$
\lean{Rm\_eq\_zero\_of\_Defm\_eq\_zero}
\end{theorem}

\begin{proof}
From the master regression theorem, if all deformation components vanish, then $2R^a_{bce} = 0$. Multiplying by $1/2$ (which exists in our ring) gives $R^a_{bce} = 0$.
\end{proof}

\begin{theorem}
Consequently Ricci vanishes as well: if $D_{ij} = 0$ for all $i,j$, then $R_{be} = 0$. \lean{Ric\_eq\_zero\_of\_Defm\_eq\_zero}
\end{theorem}

\begin{proof}
The Ricci tensor is the trace of the Riemann tensor, $R_{be} = \sum_a R^a_{bae}$. If $R^a_{bae} = 0$ for all $a$, the sum vanishes.
\end{proof}

\begin{theorem}
The deformation tensor is recovered from curvature by tracing, so $D$ carries exactly the same information as $R$:
$$2R_{be} = -(n-2)D_{be} - \delta_{be}(\text{tr}\,D),$$
where $\text{tr}\,D = 2\Delta\sigma + (n-2)|\nabla\sigma|^2$. \lean{two\_Ric\_eq}
\end{theorem}

\begin{proof}
From the Ricci formula $R_{be} = -(n-2)(\sigma_{be} - \sigma_b\sigma_e) - \delta_{be}(\Delta\sigma + (n-2)|\nabla\sigma|^2)$ and the definition of $D_{be}$, direct algebra gives the stated form.
\end{proof}

\subsection{Frame}

The original axiom (A4) postulated a single scalar scale $\sigma$, the same in every direction. This axiom is provably unable to carry Schwarzschild: if the temporal and radial scales are reciprocal and the scale is isotropic, the conformal factor must be trivial and the geometry must be flat. The diagnosis is physical: energy-momentum is a tensor, and a tensor source cannot set a scalar scale. A4 is therefore weakened to A4' (directional scale), which allows one scale per direction. Under isotropy all directions carry the same scale, recovering the old A4 as a special case; away from isotropy, genuine non-flat solutions to the reciprocal relation exist.

\begin{definition}
A4' (directional scale): One local unit of measure per direction, each invertible. The local scale in direction $a$ is $s_a$, a unit of the ring. \lean{Frame.DirScale}
\end{definition}

\begin{definition}
The measured metric component in direction $a$, with signature $\eta$: $g_a = \eta_a s_a^2$. \lean{Frame.DirScale.metric}
\end{definition}

\begin{definition}
Isotropy: all directions carry the same scale, $s_a = s_b$ for all $a,b$. \lean{Frame.DirScale.Isotropic}
\end{definition}

\begin{theorem}
Under isotropy the metric is the bare one times a single conformal factor. Everything proved in the conformal sector therefore survives as the isotropic special case of A4'. \lean{Frame.DirScale.isotropic\_metric}
\end{theorem}

\begin{proof}
If $s_a = s_b$ for all $a,b$, then all metric components $g_a = \eta_a s_a^2$ and $g_b = \eta_b s_a^2$ share the same scale factor, differing only by their signatures.
\end{proof}

\begin{definition}
The reciprocal-scale relation: $s_t \cdot s_r = 1$. This is equivalent to $g_t \cdot g_r = -1$ for signature $(\eta_t, \eta_r) = (-1, 1)$, and is the defining feature of the Schwarzschild and Reissner-Nordström families. In scale language: the temporal and radial units are exact reciprocals. \lean{Frame.DirScale.Reciprocal}
\end{definition}

\begin{theorem}
The reciprocal relation written on metric components: if $s_t \cdot s_r = 1$, $\eta_t = -1$, and $\eta_r = 1$, then $g_t \cdot g_r = -1$. \lean{Frame.DirScale.reciprocal\_metric}
\end{theorem}

\begin{proof}
We have $g_t \cdot g_r = (-1)\cdot s_t^2 \cdot (1) \cdot s_r^2 = -(s_t \cdot s_r)^2 = -(1)^2 = -1$.
\end{proof}

\begin{theorem}
Under the old, isotropic axiom the Schwarzschild relation collapses to flat space. If all directions share one scale and the temporal and radial scales are reciprocal, then that scale squares to 1: $s_t^2 = 1$. The conformal factor is trivial and the geometry is the bare Euclidean bookkeeping. So the old axiom did not merely fail to reproduce Schwarzschild---it excluded it. The axiom was too strong. \lean{Frame.DirScale.isotropic\_reciprocal\_forces\_flat}
\end{theorem}

\begin{proof}
If $s_a = s_b$ for all $a,b$, then $s_t = s_r$. The reciprocal condition $s_t \cdot s_r = 1$ becomes $s_t^2 = 1$.
\end{proof}

\begin{theorem}
Consequently every metric component is the bare one: the geometry is flat. If the isotropic and reciprocal conditions hold, then $g_a = \eta_a$ for all $a$. \lean{Frame.DirScale.isotropic\_reciprocal\_metric\_bare}
\end{theorem}

\begin{proof}
From $s_t^2 = 1$ and $s_a = s_t$ for all $a$, we have $s_a^2 = 1$. Thus $g_a = \eta_a \cdot s_a^2 = \eta_a \cdot 1 = \eta_a$.
\end{proof}

\begin{theorem}
Under A4' the Schwarzschild relation is satisfiable non-trivially. For any unit $u$ there exists a directional scale field with $s_t = u$ and $s_r = u^{-1}$, reciprocal by construction, and isotropic only in the degenerate case $u^2 = 1$. The temporal and radial scales genuinely differ, which is what Schwarzschild requires and what the scalar axiom could not provide. \lean{Frame.exists\_reciprocal\_nonisotropic}
\end{theorem}

\begin{proof}
Define a directional scale by setting $s_a = u$ if $a = t$, $s_a = u^{-1}$ if $a = r$, and $s_a = 1$ otherwise. Then $s_t \cdot s_r = u \cdot u^{-1} = 1$ by the unit axiom. For isotropy to hold, we would need $u = u^{-1}$, which implies $u^2 = 1$, contradicting the hypothesis $u^2 \neq 1$.
\end{proof}

\subsection{Unify}

The directional scale A4' is the axiom; A4 (scalar scale) is a theorem about its isotropic locus. This unifies a potentially inconsistent pair of statements: the conformal sector developed under A4, and the Frame sector under A4', now coexist as the isotropic sector of the latter. Every theorem previously proved about $\sigma$ holds verbatim, now with a sharp domain of validity. Nothing is lost and the axiom list has one entry where it had two.

\begin{definition}
The measured metric of a directional scale field: diagonal in the frame, with the local unit $s_a$ in direction $a$ and signature $\eta$:
$$g_{ab} = \delta_{ab} \eta_a s_a^2.$$
\lean{Unify.toMetric}
\end{definition}

\begin{definition}
The metric of the old axiom A4: one scalar scale, $g_{ab} = \delta_{ab}\eta_a u^2$ where $u$ is the scale. \lean{Unify.conformalMetric}
\end{definition}

\begin{theorem}
The old axiom is the isotropic locus of the new one. If all directions carry the same scale $s$, the metric of A4' is literally the conformal metric of A4. Every result proved in the conformal sector therefore holds unchanged, as a statement about isotropic scale fields. \lean{Unify.toMetric\_of\_isotropic}
\end{theorem}

\begin{proof}
If $s_a = s_c$ for all $a,c$, then $g_{ab} = \delta_{ab}\eta_a s_a^2 = \delta_{ab}\eta_a s_c^2$, which is the conformal metric.
\end{proof}

\begin{definition}
Embedding of a scalar scale: every scalar scale is a directional scale, namely an isotropic one. Specifically, $s_a = u$ for all $a$. \lean{Unify.ofScalar}
\end{definition}

\begin{theorem}
The scalar embedding is isotropic: the scale $\text{ofScalar}(u)$ has $s_a = u$ for all $a$, so $s_a = s_b$ for all pairs $(a,b)$. \lean{Unify.ofScalar\_isotropic}
\end{theorem}

\begin{proof}
Immediate from the definition.
\end{proof}

\begin{theorem}
The metric of an embedded scalar scale is the conformal metric:
$$\text{ofScalar}(u).toMetric(\eta) = conformalMetric(u, \eta).$$
\lean{Unify.ofScalar\_toMetric}
\end{theorem}

\begin{proof}
Since $s_a = u$ for all $a$, we have $g_{ab} = \delta_{ab}\eta_a u^2$, which is the conformal metric.
\end{proof}

\begin{theorem}
The isotropic locus is exactly the image of the scalar theory. An $n$-direction scale field is isotropic if and only if it is $\text{ofScalar}(u)$ for some unit $u$. So the old axiom is not merely implied by the new one---it is precisely the new one restricted to a sub-locus, with nothing left over. \lean{Unify.isotropic\_iff\_ofScalar}
\end{theorem}

\begin{proof}
($\Rightarrow$) If the scale is isotropic, all directions carry $s_a = s_{\langle 0 \rangle}$ for all $a$, so the scale field is $\text{ofScalar}(s_{\langle 0 \rangle})$. ($\Leftarrow$) If the scale is $\text{ofScalar}(u)$, then $s_a = u$ for all $a$, so it is isotropic.
\end{proof}

Scalar scale fields (A2/A3) embed as isotropic directional scales. The metric induced by the directional geometry is identical to the conformal metric of the scalar theory.

\begin{definition}
A scalar scale field $F$, viewed as an isotropic directional scale field. \lean{ScaleField.toDirScale}
\end{definition}

\begin{theorem}
The embedding is isotropic: $F.toDirScale.Isotropic$ for any scalar field $F$. \lean{Unify.toDirScale\_isotropic}
\end{theorem}

\begin{proof}
Immediate from $\text{ofScalar\_isotropic}$.
\end{proof}

\begin{theorem}
The conformal sector sits inside A4' exactly. The metric obtained from A2/A3 by the new axiom is the metric the old axiom postulated:
$$F.toDirScale.toMetric(\eta) = conformalMetric(F.s, \eta).$$
Consequently the curvature computed in Conformal.lean from $F.\sigma$ is the curvature of the A4' metric of $F.toDirScale$. \lean{Unify.toDirScale\_toMetric}
\end{theorem}

\begin{proof}
Apply $\text{ofScalar\_toMetric}$ to $F.s$.
\end{proof}

\begin{theorem}
Collected coherence statement. For any scalar scale field $F$:
\begin{enumerate}
\item its directional embedding is isotropic;
\item the metric it induces under A4' equals the conformal metric of A4;
\item consequently the curvature computed in Conformal.lean from $F.\sigma$ is the curvature of the A4' metric of $F.toDirScale$.
\end{enumerate}
This is the precise sense in which the earlier development is a sub-theory of the present one rather than a competitor to it. \lean{Unify.sector\_coherence}
\end{theorem}

\begin{proof}
Combine the three preceding theorems.
\end{proof}

The two sectors are distinguished by a single observable: whether the metric is a multiple of the bare bookkeeping.

\begin{theorem}
Under isotropy the metric is a common multiple of the signature. If $s_a = s_b$ for all $a,b$, then
$$g_{aa}\eta_b = g_{bb}\eta_a.$$
\lean{Unify.isotropic\_ratio\_const}
\end{theorem}

\begin{proof}
$g_{aa} = \eta_a s_a^2$ and $g_{bb} = \eta_b s_a^2$, so $g_{aa}\eta_b = \eta_a s_a^2 \cdot \eta_b = \eta_a \cdot \eta_b s_a^2 = g_{bb}\eta_a$.
\end{proof}

\begin{theorem}
Conversely, if two directions carry different squared scales the field is not isotropic. If $s_a^2 \neq s_b^2$, then the scale field is not isotropic. This is the observable that Schwarzschild needs and that the scalar axiom could not supply. \lean{Unify.not\_isotropic\_of\_ne}
\end{theorem}

\begin{proof}
If the scale is isotropic, then $s_a = s_b$, so $s_a^2 = s_b^2$. Contrapositive: if $s_a^2 \neq s_b^2$, the scale is not isotropic.
\end{proof}

\subsection{Newton}

Given the regression theorem that relates gravity to scale field inhomogeneity, one further statement---a scale response law---closes the system. The statement is that the scalar curvature of the measured metric is proportional to the local energy density. In the first-order (weak-field) regime, this constraint is exactly Poisson's equation, with the log-scale identified as minus the Newtonian potential. Both signs emerge naturally without being imposed. The first-order regime is not merely an approximation; it is realized exactly by extending the substrate to dual numbers $A[\epsilon]/(\epsilon^2)$.

\begin{definition}
The Newtonian potential is minus the log-scale: $\Phi = -\sigma$. Deep in a gravitational well the scale is enlarged ($\sigma > 0$), so $\Phi < 0$ as it must be. \lean{Newton.newtPot}
\end{definition}

\begin{theorem}
The Laplacian of the Newtonian potential: $\Delta\Phi = -\Delta\sigma$. \lean{Newton.lap\_newtPot}
\end{theorem}

\begin{proof}
Immediate from $\Phi = -\sigma$ and linearity of the Laplacian.
\end{proof}

\begin{theorem}
In the first-order regime the scalar curvature is linear in the scale: the $|\nabla\sigma|^2$ term drops out identically. If $|\nabla\sigma|^2 = 0$, then
$$R_{\text{bare}} = -2(n-1)\Delta\sigma.$$
\lean{Newton.RscBare\_linear}
\end{theorem}

\begin{proof}
From Conformal.RscBare\_eq, drop the term $-(n-1)(n-2)|\nabla\sigma|^2$ when $|\nabla\sigma|^2 = 0$.
\end{proof}

\begin{theorem}
Poisson's equation. From the scale response law $e^{2\sigma}R = \kappa\rho$ in the first-order regime,
$$2(n-1)\Delta\Phi = \kappa\rho, \quad \Phi = -\sigma.$$
For $n = 3$ and $\kappa = 16\pi G$ this is $\Delta\Phi = 4\pi G\rho$ exactly. Newtonian gravity is therefore not an input: it is the linearization of the statement that energy sets the local scale. \lean{Newton.poisson}
\end{theorem}

\begin{proof}
From the scale response law and the linear regime, $R_{\text{bare}} = -2(n-1)\Delta\sigma = \kappa\rho$. Substituting $\sigma = -\Phi$ gives $2(n-1)\Delta\Phi = \kappa\rho$.
\end{proof}

\begin{theorem}
The physical case $n = 3$: $4\Delta\Phi = \kappa\rho$, i.e. $\Delta\Phi = 4\pi G\rho$ for $\kappa = 16\pi G$. \lean{Newton.poisson\_three}
\end{theorem}

\begin{proof}
From Poisson's equation with $n = 3$, we have $2(3-1)\Delta\Phi = \kappa\rho$, i.e., $4\Delta\Phi = \kappa\rho$.
\end{proof}

The first-order regime is not a hypothesis we hope is harmless. Extend the substrate by a square-zero infinitesimal, $A[\epsilon]/(\epsilon^2)$, and let the scale be infinitesimal, $\sigma = \epsilon\varphi$. Then $|\nabla\sigma|^2 = \epsilon^2|\nabla\varphi|^2 = 0$ identically, and all results above hold with no approximation whatsoever. Poisson's equation is exact.

\begin{theorem}
The derivations extend componentwise to the dual numbers $A[\epsilon]/(\epsilon^2)$: $\partial_i(x, y) = (\partial_i x, \partial_i y)$ where the first component is the ordinary derivative and the second is the dual part. \lean{Newton.Dual.instScaleAlgebra}
\end{theorem}

\begin{proof}
The axioms of $ScaleAlgebra$ (additivity and Leibniz rule of derivations, commutativity of mixed partials) are verified componentwise.
\end{proof}

\begin{theorem}
An infinitesimal scale field $\sigma = \epsilon\varphi$ has purely infinitesimal gradient: $\sigma_i = \epsilon(\partial_i\varphi)$. \lean{Newton.Dual.sig\_inr}
\end{theorem}

\begin{proof}
$\sigma_i = \partial_i \sigma = \partial_i(\epsilon\varphi) = \epsilon\partial_i\varphi$ by the Leibniz rule applied componentwise.
\end{proof}

\begin{theorem}
The linearised regime is exact. For an infinitesimal scale field the quadratic invariant vanishes identically---no term is discarded:
$$|\nabla\sigma|^2 = 0 \quad \text{for } \sigma = \epsilon\varphi.$$
\lean{Newton.Dual.gradsq\_inr}
\end{theorem}

\begin{proof}
$|\nabla\sigma|^2 = \sum_i \sigma_i^2 = \sum_i (\epsilon\partial_i\varphi)^2 = \sum_i \epsilon^2(\partial_i\varphi)^2 = 0$ since $\epsilon^2 = 0$ in the dual numbers.
\end{proof}

\begin{theorem}
Consequently Poisson's equation holds exactly for an infinitesimal scale field: no linearization error, because there is nothing to linearize.
$$2(n-1)\Delta\Phi = \kappa\alpha$$
for $\Phi = -\epsilon\varphi$ and $\alpha$ the energy density in dual numbers. \lean{Newton.Dual.poisson\_exact}
\end{theorem}

\begin{proof}
Apply the Poisson equation theorem with $\sigma = \epsilon\varphi$ and the exactness of $|\nabla\sigma|^2 = 0$.
\end{proof}

\subsection{Schwarzschild}

The vacuum condition (vanishing Ricci) does not directly imply reciprocity $s_t \cdot s_r = 1$. It implies something weaker: the sum of temporal and radial log-scales is constant. Getting from constant to one is a separate step, and it is the step where physical content appears---it is asymptotic flatness, the requirement that far from the source the local unit agrees with the fiducial. But deriving the constancy itself requires computing the Ricci tensor for the diagonal ansatz, which is the outstanding piece of the gravity chain.

This file proves the part of that computation where the non-obvious cancellation happens. Write the metric in directional form with only temporal and radial scales varying, $A = e^{2\sigma_t}$, $B = e^{2\sigma_r}$. The two Ricci components are complicated expressions in $A, B$ and their derivatives; forming the particular combination $R_t^t/A + R_r^r/B$ makes all the awkward terms cancel, leaving only a total logarithmic derivative of the product $AB$.

\begin{definition}
The temporal component of the Ricci tensor for $-A\,dt^2 + B\,dr^2 + r^2 d\Omega^2$, where primes denote $r$-derivatives:
$$R_{tt} = \frac{A''}{2B} - \frac{A'}{4B}\left(\frac{A'}{A} + \frac{B'}{B}\right) + \frac{A'}{rB}.$$
\lean{Schwarzschild.Rtt}
\end{definition}

\begin{definition}
The radial component:
$$R_{rr} = -\frac{A''}{2A} + \frac{A'}{4A}\left(\frac{A'}{A} + \frac{B'}{B}\right) + \frac{B'}{rB}.$$
\lean{Schwarzschild.Rrr}
\end{definition}

\begin{theorem}
Everything awkward cancels. The particular combination
$$\frac{R_{tt}}{A} + \frac{R_{rr}}{B} = \frac{1}{rB}\left(\frac{A'}{A} + \frac{B'}{B}\right).$$
The second-order derivatives cancel against each other and so do the quadratic terms, exactly and for all $A, B$. What survives is a total logarithmic derivative---which is why the vacuum constrains the product $AB$ and nothing else about it. \lean{Schwarzschild.ricci\_combination}
\end{theorem}

\begin{proof}
Substitute the definitions of $R_{tt}$ and $R_{rr}$, clear denominators using $A \neq 0$, $B \neq 0$, $r \neq 0$, and collect terms by degree and structure. All second-derivative and quadratic terms cancel, leaving only the logarithmic derivative term.
\end{proof}

\begin{theorem}
In vacuum the logarithmic derivative of the product vanishes. If $R_{tt} = 0$ and $R_{rr} = 0$, then
$$\frac{A'}{A} + \frac{B'}{B} = 0,$$
because the prefactor $1/(rB)$ is never zero. \lean{Schwarzschild.vacuum\_log\_derivative}
\end{theorem}

\begin{proof}
From the cancellation theorem, if both Ricci components vanish, then $\frac{1}{rB}(\frac{A'}{A} + \frac{B'}{B}) = 0$. Since $r \neq 0$ and $B \neq 0$, the prefactor is nonzero, so the logarithmic derivative must vanish.
\end{proof}

\begin{theorem}
In scale variables, the vacuum condition reads as the scale sum being stationary. With $A' / A = 2\sigma_t'$ and $B' / B = 2\sigma_r'$, the condition $A'/A + B'/B = 0$ becomes
$$\sigma_t' + \sigma_r' = 0,$$
which is exactly the hypothesis needed in Vacuum.reciprocity\_of\_asymptotic\_flatness. \lean{Schwarzschild.scale\_sum\_derivative\_zero}
\end{theorem}

\begin{proof}
Substitute the scale-variable substitutions into the logarithmic derivative condition: $2\sigma_t' + 2\sigma_r' = 0$ gives $\sigma_t' + \sigma_r' = 0$.
\end{proof}

\begin{theorem}
The vacuum condition on the diagonal ansatz is exactly the statement that the scale sum is stationary. Nothing else about $A$ and $B$ is constrained by it, which is why the constant value of the scale sum must come from a boundary condition. \lean{Schwarzschild.vacuum\_iff\_scale\_sum\_stationary}
\end{theorem}

\begin{proof}
Combine the three preceding results.
\end{proof}

\begin{theorem}
The combination that cancels is the one weighted by the inverse metric on the $t$--$r$ block. The area element in that plane is $AB$, and its logarithmic derivative is what the vacuum sets to zero:
$$A' B + A B' = 0.$$
Read physically: in vacuum, gravity trades temporal scale against radial scale and creates no scale in that plane at all. \lean{Schwarzschild.area\_element\_stationary}
\end{theorem}

\begin{proof}
From the vacuum condition $A'/A + B'/B = 0$, clear the common denominator to get $A'B + AB' = 0$.
\end{proof}

\subsection{RicciDiag}

The Schwarzschild.lean file verified the crucial cancellation that makes reciprocity work, but treated the two Ricci components as given. This file derives them from the metric using the Levi-Civita connection, closing the loop. For a diagonal metric $g_{tt} = -A(r)$, $g_{rr} = B(r)$, $g_{\theta\theta} = r^2$, the connection coefficients follow from the formula $\Gamma^a_{bc} = \frac{1}{2}g^{aa}(\partial_b g_{ac} + \partial_c g_{ab} - \partial_a g_{bc})$. The Ricci components then come from summing the connection and its derivatives according to the definition. The point of this file is that the resulting expressions are exactly those Schwarzschild.lean assumed.

\begin{definition}
Connection coefficients. For the diagonal metric with $g_{tt} = -A$, $g_{rr} = B$, $g_{\theta\theta} = r^2$:
$$\Gamma^t_{tr} = \frac{A'}{2A}.$$
\lean{RicciDiag.Γ\_t\_tr}
\end{definition}

\begin{definition}
$$\Gamma^r_{tt} = \frac{A'}{2B}.$$
\lean{RicciDiag.Γ\_r\_tt}
\end{definition}

\begin{definition}
$$\Gamma^r_{rr} = \frac{B'}{2B}.$$
\lean{RicciDiag.Γ\_r\_rr}
\end{definition}

\begin{definition}
$$\Gamma^\theta_{r\theta} = \Gamma^\phi_{r\phi} = \frac{1}{r}.$$
\lean{RicciDiag.Γ\_ang\_r}
\end{definition}

\begin{definition}
The contracted symbol:
$$\Gamma^c_{cr} = \Gamma^t_{tr} + \Gamma^r_{rr} + 2\Gamma^\theta_{r\theta}.$$
\lean{RicciDiag.Γtrace\_r}
\end{definition}

\begin{definition}
Radial derivatives of the connection coefficients. By the quotient rule:
$$\partial_r \Gamma^r_{tt} = \frac{A''}{2B} - \frac{A'B'}{2B^2}.$$
\lean{RicciDiag.dΓ\_r\_tt}
\end{definition}

\begin{definition}
$$\partial_r \Gamma^t_{tr} = \frac{A''}{2A} - \frac{(A')^2}{2A^2}.$$
\lean{RicciDiag.dΓ\_t\_tr}
\end{definition}

\begin{definition}
$$\partial_r \Gamma^\theta_{r\theta} = -\frac{1}{r^2}.$$
\lean{RicciDiag.dΓ\_ang\_r}
\end{definition}

\begin{definition}
The temporal Ricci component $R_{tt} = \partial_r \Gamma^r_{tt} + \Gamma^c_{cr}\Gamma^r_{tt} - 2\Gamma^t_{tr}\Gamma^r_{tt}$, using the fact that the metric is static so no $\partial_t$ survives and $\Gamma^c_{tt}$ is nonzero only for $c = r$:
$$R^{\text{Ricci}}_{tt} = \partial_r \Gamma^r_{tt} + \Gamma^c_{cr}\Gamma^r_{tt} - 2\Gamma^t_{tr}\Gamma^r_{tt}.$$
\lean{RicciDiag.ricciTT}
\end{definition}

\begin{definition}
The radial Ricci component:
$$R^{\text{Ricci}}_{rr} = -\partial_r \Gamma^t_{tr} - 2\partial_r \Gamma^\theta_{r\theta} + \Gamma^c_{cr}\Gamma^r_{rr} - (\Gamma^t_{tr})^2 - (\Gamma^r_{rr})^2 - 2(\Gamma^\theta_{r\theta})^2.$$
\lean{RicciDiag.ricciRR}
\end{definition}

\begin{theorem}
The temporal component, derived. Evaluating the Ricci sum reproduces exactly the expression Schwarzschild.Rtt assumed:
$$R^{\text{Ricci}}_{tt} = R_{tt}.$$
\lean{RicciDiag.ricciTT\_eq}
\end{theorem}

\begin{proof}
Substitute the definitions of all connection coefficients and their derivatives, clear denominators, and simplify using ring identities. The result matches the Schwarzschild.Rtt expression.
\end{proof}

\begin{theorem}
The radial component, derived:
$$R^{\text{Ricci}}_{rr} = R_{rr}.$$
\lean{RicciDiag.ricciRR\_eq}
\end{theorem}

\begin{proof}
Substitute the definitions and simplify as above.
\end{proof}

\begin{theorem}
The gravity chain, end to end. From the connection of the diagonal metric: the Ricci components are what Schwarzschild.lean assumed, their particular combination cancels, and the vacuum therefore makes the scale sum stationary.

Given $A = e^{2\sigma_t}$, $B = e^{2\sigma_r}$ with $A' / A = 2\sigma_t'$ and $B' / B = 2\sigma_r'$, and assuming the vacuum condition $R^{\text{Ricci}}_{tt} = 0$ and $R^{\text{Ricci}}_{rr} = 0$, we have
$$\sigma_t' + \sigma_r' = 0.$$

With Vacuum.reciprocity\_of\_asymptotic\_flatness this gives $s_t \cdot s_r = 1$. No step between the metric and that conclusion is now an assumption. \lean{RicciDiag.gravity\_chain}
\end{theorem}

\begin{proof}
From ricciTT\_eq and ricciRR\_eq, rewrite the Ricci conditions. Then apply Schwarzschild.vacuum\_iff\_scale\_sum\_stationary to conclude.
\end{proof}

\begin{theorem}
Consequently the $t$--$r$ area element is stationary in vacuum, derived now from the metric rather than assumed. If the Ricci components vanish, then $A'B + AB' = 0$. \lean{RicciDiag.area\_stationary\_derived}
\end{theorem}

\begin{proof}
Rewrite the Ricci conditions using ricciTT\_eq and ricciRR\_eq, then apply Schwarzschild.area\_element\_stationary.
\end{proof}

\subsection{Vacuum}

The vacuum condition in the standard static spherically symmetric computation does not directly give $s_t s_r = 1$. It gives something weaker: the sum of temporal and radial log-scales has vanishing derivative, so $s_t s_r$ is constant. Getting from constant to one requires an additional boundary condition: asymptotic flatness, the physical requirement that far from the source there is nothing, so the local unit must agree with the fiducial. This file derives reciprocity from the vacuum consequence plus asymptotic flatness.

\begin{theorem}
A function with vanishing derivative that vanishes at one point vanishes everywhere. If $f'(r) = 0$ for all $r$ and $f(r_0) = 0$ for some $r_0$, then $f(r) = 0$ for all $r$. \lean{Vacuum.const\_of\_deriv\_zero\_and\_vanishing}
\end{theorem}

\begin{proof}
A function with zero derivative everywhere is constant. If it equals zero at one point, the constant must be zero, so the function vanishes everywhere.
\end{proof}

\begin{theorem}
Reciprocity from asymptotic flatness. The vacuum condition makes the sum of the temporal and radial log-scales constant. Requiring that the sum vanish far from the source---asymptotic flatness, i.e. that the local unit agrees with the fiducial where there is no matter---forces the sum to vanish everywhere:
$$\forall r,\quad \sigma_t(r) + \sigma_r(r) = 0.$$

The constant is therefore not fitted; it is fixed by the only boundary condition the axioms permit, since by A5 the fiducial is exactly what ``far away'' means. \lean{Vacuum.reciprocity\_of\_asymptotic\_flatness}
\end{theorem}

\begin{proof}
Apply const\_of\_deriv\_zero\_and\_vanishing to the sum $f(r) = \sigma_t(r) + \sigma_r(r)$, using the hypothesis that $f'(r) = 0$ (from the vacuum condition) and $f(r_\infty) = 0$ (from asymptotic flatness).
\end{proof}

\begin{theorem}
The same statement multiplicatively: the temporal and radial scales are exact reciprocals,
$$s_t(r) \cdot s_r(r) = e^{\sigma_t(r)} e^{\sigma_r(r)} = e^{\sigma_t(r) + \sigma_r(r)} = 1.$$
\lean{Vacuum.scale\_product\_eq\_one}
\end{theorem}

\begin{proof}
From reciprocity\_of\_asymptotic\_flatness, $\sigma_t(r) + \sigma_r(r) = 0$, so $e^{\sigma_t(r) + \sigma_r(r)} = e^0 = 1$.
\end{proof}

\begin{theorem}
Without the boundary condition, reciprocity fails: the scale product is some constant, and any constant is compatible with the vacuum equation. The product is constant:
$$s_t(r_1) s_r(r_1) = s_t(r_2) s_r(r_2)$$
for any two points $r_1, r_2$. Asymptotic flatness is therefore doing real work and is not a convention. \lean{Vacuum.product\_constant\_without\_flatness}
\end{theorem}

\begin{proof}
From the vacuum condition that the sum $\sigma_t + \sigma_r$ has zero derivative, it is constant. Hence $e^{\sigma_t(r_1) + \sigma_r(r_1)} = e^{\sigma_t(r_2) + \sigma_r(r_2)}$, which gives the product equality.
\end{proof}

\begin{theorem}
The area element in the time--radial plane is what reciprocity conserves. The t--r area element is $AB = e^{2\sigma_t + 2\sigma_r}$. Under reciprocity (which sets $\sigma_t + \sigma_r = 0$ everywhere), we have
$$e^{2(\sigma_t + \sigma_r)} = e^0 = 1.$$
Read physically: in vacuum, gravity trades temporal scale against radial scale and creates no scale in that plane at all. \lean{Vacuum.tr\_area\_element\_constant}
\end{theorem}

\begin{proof}
From reciprocity\_of\_asymptotic\_flatness, $\sigma_t(r) + \sigma_r(r) = 0$, so $e^{2(\sigma_t + \sigma_r)} = e^0 = 1$.
\end{proof}

